
\section{Сравнительный анализ подходов}

%% а́~— я́ о́~— ё э́~— е́ у́~— ю́ ы́~— и́

Важно понимать разницу между~поиском нечетких дубликатов и~их идентификацией.

\subsection{Поиск}

Поиск применяется, когда существует база данных видео.
Пользователи вводят видео-запрос через~поисковый интерфейс.
Поисковая система получает запрос, выделяет его характеристики и~сигнатуры,
и~сравнивает их с~данными из~базы.
Пользователям возвращается список результатов,
в~порядке убывания похожести или~в~порядке возрастания «расстояния» между
между~видео из~базы и~из~запроса.
Отношения между~дубликатами можно представить и~матрицей истинности.
Но~отсортированный список может правильно
описать разницу между~их результатам.

Бо́льшая часть существующих исследований посвящены задаче
поиска \cite{Liu:2013}.

\subsection{Идентификация}

Существует два~случая идентификации нечетких дубликатов.

В~первом случае, идентификация направлена на~формирование групп НДВ
из~одного или~нескольких заданных множеств видео
\cite{Tan:2008} \cite{Wu:2009} \cite{Zhao:2009}.
Так~решается проблема фильтрации баз~видео,
определения жанра, защиты авторских прав, и~т.~д.
Поиск и~этот вариант идентификации совместно используют
общих методов и~процедур, особенно для~извлечения признаков из~видео,
генерации и~сопоставления сигнатур видео.
Однако, при~идентификации приходится перебирать большое число вариантов,
что~увеличивает время затраченное на~получение результата,
по~сравнению с~поиском.

Второй случай~— обнаружение дубликатов в~реальном времени,
где запрос передается в~форме непрерывного видео-потока.
Такой сценарий предполагает, что, в~отличие от~поиска, система не~получает
ни~полную информацию о~видео, ни~само видео целиком.
Записать целое видео невозможно,
так как~сам поток может быть потенциально бесконечным, например,
в~случае камеры наружного наблюдения или~телевизионного эфира.
С~точки зрения системы обнаружения, на~каждом конкретном промежутке времени,
данные будут неполными. Другая проблема состоит в~том, что~из-за~непрерывного
появления новых данных в~потоке, ранее полученные данные должны
быть обработаны своевременно.
В~результате этих ограничений методы обнаружения дубликатов,
которые используются при~поиске не~могут быть использованы.

По~сравнению с~поиском и~первым случаем идентификации,
существует не~много работ, посвященных обнаружению дубликатов
в~реальном времени: \cite{Xie:2010}, \cite{Huang:2009},
\cite{Shen:2007}, \cite{Yan:2008}. В~некоторых приложениях требуется найти
точное место появления нечеткого дубликата \cite{Liu:2013}.

